\documentclass[11pt,a4paper]{article}
\usepackage[margin=2.3cm]{geometry}
\usepackage{amsmath,amssymb}
\usepackage{booktabs,longtable,array}
\usepackage[version=4]{mhchem}
\usepackage{xcolor}
\usepackage[colorlinks=true,linkcolor=blue!60!black,urlcolor=blue!60!black]{hyperref}
\usepackage{enumitem}
\usepackage{fancyhdr}
\pagestyle{fancy}\fancyhf{}
\rhead{GreenhouseTwin n79}\lhead{Model description}\cfoot{\thepage}
\newcommand{\code}[1]{\texttt{#1}}
\sloppy

\title{\textbf{GreenhouseTwin n79}\\[4pt]\large The greenhouse--cucumber model without flags:\\
climate, photosynthesis, crop growth and LAI}
\author{Greenhouses project --- model reference}
\date{Branch \code{n79}, 21 October 2026}

\begin{document}
\maketitle
\begin{abstract}
\noindent This document specifies the model coded in \code{invernaderosN79/GreenhouseTwin}: one equation set, no
behaviour flags. It is the configuration chosen as the baseline for making the physics more descriptive: the V2.2
engine with the n79 settings (fixed compartment geometry, energy-conserving solar column, cover node with
condensation, open-loop logged vents, stack leakage, closed floor balance, internal mass, full lamp heat, crop
respiration) and Aaron V\'elez's generic cucumber model with its canopy and LAI extension (the n105 crop set-up).
Every equation is written with the term names of the code (\code{src/climate.jl}, \code{src/crop.jl}), so each line can
be checked against \code{model\_terms}. Parameter tables are generated from the parameter files. Known structural
gaps are listed at the end; they are the starting point of the physics work.
\end{abstract}
\tableofcontents

%======================================================================
\section{Overview}\label{sec:overview}
%======================================================================
\paragraph{Scope.} One compartment of the Autonomous Greenhouse Challenge 2018 (AiCU, Bleiswijk; 96\,m$^2$, Venlo, gutter
5\,m), cucumber crop transplanted on 14 August 2018, 115 simulated days. All quantities are per m$^2$ of floor.

\paragraph{Two models, offline coupling.}
\begin{itemize}[nosep]
\item The \textbf{climate model} is a system of 8 ODEs integrated day by day (Tsit5, 5-min output):
\begin{equation}
u=\bigl[\,T_1,\;T_2,\;C_1,\;V_1,\;e_{\rm int},\;T_5,\;T_d,\;T_m\,\bigr],
\end{equation}
canopy temperature, air temperature, \ce{CO2} concentration [mg\,m$^{-3}$], vapour pressure [Pa], integral of the heating error,
floor layer, deep soil and internal mass temperatures [K].
\item The \textbf{crop model} is a daily step (Section~\ref{sec:crop}) driven by measured climate. Its daily LAI, non-leaf
respiration and standing fresh mass are inputs of the climate model (\code{test/crop\_series.jl}). The crop does not see
the simulated climate (live coupling is a planned extension).
\end{itemize}

\paragraph{Provenance and verification.} The code is GreenhouseSim V2.2 (git tag \code{v2.2-flags-final}) with every flag
resolved to its n79 value and dead branches removed. \code{test/equivalence\_synthetic.jl} and
\code{test/equivalence\_n79.jl} compare it with V2.2 run with the n79 flags (Section~\ref{sec:verify}).

\paragraph{Inputs.} Table~\ref{tab:inputs}. Weather and LAI are functions of absolute time; controls and crop inputs of
time since the season start. All are 5-min step functions.
\begin{table}[h]\centering\small
\caption{Model inputs (\code{src/inputs.jl}; AGC-2018 construction in \code{src/agc\_data.jl}).}\label{tab:inputs}
\begin{tabular}{>{\raggedright\arraybackslash}p{0.17\textwidth}>{\raggedright\arraybackslash}p{0.24\textwidth}>{\raggedright\arraybackslash}p{0.5\textwidth}}
\toprule symbol & code & meaning (AiCU source)\\\midrule
$I_{\rm glob}$ & \code{Iglob} & global radiation outside [W\,m$^{-2}$] (meteo)\\
$T_{\rm sky}$ & \code{Tsky} & sky temperature from the pyrgeometer [K]\\
$T_o$, $v$, $V_o$ & \code{Tout, Wind, VPout} & outside temperature, wind speed, vapour pressure\\
$C_o$, $T_{\rm mc}$ & \code{CO2out, TmechCool} & outside \ce{CO2} 834.7\,mg\,m$^{-3}$; mechanical-cooling surface 273.15\,K\\
$L$ & \code{LAI} & leaf area index per m$^2$ floor (crop series)\\
$U_1$, $U_{\rm bl}$ & \code{EnScr, BlackScr} & energy and blackout screen closure\\
$U_{\rm lee}$, $U_{\rm win}$ & \code{VentLee, VentWind} & leeward, windward roof-vent openings\\
$U_{10}$ & \code{CO2Valve} & \ce{CO2} valve opening (setpoint rule scaled to the logged daily dosage)\\
$U_{11}$ & \code{HeatSP} & heating setpoint [K]\\
$U_{12}$ & \code{Lamps} & lamp fraction\\
$T_{\rm log}$ & \code{Tpipe} & logged pipe temperature [K] (both circuits folded; no reading $\to T_o+2$)\\
$R_{\rm crop}$, FW & \code{CropResp, CropFW} & crop respiration at 20\,$^\circ$C [mg\,\ce{CO2}\,m$^{-2}$\,s$^{-1}$]; standing fresh mass [kg\,m$^{-2}$]\\
$U_{2\ldots7}$, $U_9$ & \code{FanPad}, \code{MechCool}, \code{Heater}, \code{Shade}, \code{SideVent}, \code{ForcedVent}, \code{Fog} & other actuators; zero in the AiCU season\\
\bottomrule\end{tabular}
\end{table}

%======================================================================
\section{Climate model}\label{sec:climate}
%======================================================================
\subsection{State equations}
With $c_{\rm air}=\phi_2\rho_3\alpha_5$ the air heat capacity per m$^2$ floor and $m_{\rm air}$ = \code{air\_cap\_mult},
\begin{align}
\alpha_1 L\,\dot T_1 &= r_1 + Q_{lc} + r_6 - h_1 - l_1 - r_7 - r_{14},\label{eq:T1}\\
m_{\rm air}\,c_{\rm air}\,\dot T_2 &= h_1 + h_2 + h_4 - h_7 - h_{11} - r_{12} - r_{13} - l_2 + Q_{sa} - Q_{af} - Q_{am} + Q_{la},\label{eq:T2}\\
\phi_2\,\dot C_1 &= o_1 + o_2 + o_3 - A - o_5 + o_R,\label{eq:C1}\\
\kappa_3\,\dot V_1 &= p_1 + p_2 + p_3 + p_4 - p_5 - p_6 - p_7 - p_8 - p_9 - p_{10},\qquad \kappa_3=\frac{\psi_1\phi_2}{\omega_2T_2},\label{eq:V1}\\
\dot e_{\rm int} &= \text{(anti-windup integral of } U_{11}-T_2\text{, eq.~\ref{eq:pi})},\\
C_{\rm soil}\,\dot T_5 &= Q_{af} - Q_{fd} + h_{11} + r_{14} - r_{11} + Q_{fs},\label{eq:T5}\\
\dot T_d &= (T_5-T_d)/\tau_d,\\
\bigl(C_m + c_{p,\rm crop}\,\mathrm{FW}\bigr)\,\dot T_m &= Q_{am} + Q_{ms}.\label{eq:Tm}
\end{align}
The cover temperature $T_c$ and the screen temperature $T_{\rm scr}=\tfrac12(T_o+T_2)$ are algebraic. Every term is defined
below with the name it has in \code{model\_terms}.

\subsection{Solar radiation: the energy-conserving column}\label{sec:solar}
For each band $b$ (PAR $I_P=\eta_2I_{\rm glob}$, NIR $I_N=\eta_3I_{\rm glob}$, $\eta_2+\eta_3=1$) five elements are stacked
from the sky down: roof, energy screen, blackout screen, canopy, opaque floor. Each layer is symmetric with transmission
$t$ and reflection $r$:
\begin{align}
\text{roof:}\quad & t_r=(1-\eta_1)\tau_1\ (=0.65),\qquad r_r=1-t_r-\eta_1-a_c\ (=0.15),\\
\text{screen at closure }U:\quad & t=1-U(1-\tau_s),\qquad r=U\rho_s\qquad\text{(Vanthoor 2011, 8.14--8.17)},\\
\text{canopy:}\quad & t_c=e^{-kL},\qquad r_c=\rho_\infty(1-t_c)\qquad\text{(Vanthoor 8.31--8.32)},
\end{align}
with $k=\beta_1$, $\rho_\infty=\rho_1$ (PAR) and $k=\beta_3$, $\rho_\infty=$ \code{rho\_can\_nir} (NIR); floor reflectance
$\rho_2$ (PAR), $1-\alpha_7$ (NIR); screens: energy $(\tau_s,\rho_s)=(0.75,0.15)$, blackout $(0.02,0.50)$, NIR = PAR
(provisional). With all repeated reflections (adding method), from the bottom up,
\begin{align}
R_4&=r_4+\frac{t_4^2\rho_f}{1-r_4\rho_f},\quad R_j=r_j+\frac{t_j^2R_{j+1}}{1-r_jR_{j+1}}\ (j=3,2,1),\\
D_2&=\frac{t_1}{1-r_1R_2},\quad D_{j+1}=\frac{t_jD_j}{1-r_jR_{j+1}}\ (j=2,3),\quad D_5=\frac{t_4D_4}{1-r_4\rho_f},\\
A_1&=(1-t_1-r_1)(1+R_2D_2),\quad A_j=(1-t_j-r_j)(D_j+R_{j+1}D_{j+1}),\quad A_4=(1-t_4-r_4)(D_4+\rho_fD_5),\\
A_f&=(1-\rho_f)D_5,\qquad R_1+A_1+A_2+A_3+A_4+A_f=1 .
\end{align}
Booking [W\,m$^{-2}$]: canopy $r_1=Q_{\rm can}=I_PA_4^P+I_NA_4^N$; floor level $Q_{\rm flr}=I_PA_f^P+I_NA_f^N$, of which the
floor gets $Q_{fs}=(1-f_m)Q_{\rm flr}$ and the internal mass $Q_{ms}=f_mQ_{\rm flr}$; roof absorption
$Q_{\rm roof}=\sum_bI_bA_1^b$ splits into structure $Q_{st}=\sum_bI_bA_1^b\min(\eta_1/(1-t_r-r_r),1)$ (to the air) and glass (to the
cover); screen absorption $Q_{\rm scr}=\sum_bI_b(A_2^b+A_3^b)$ goes \code{scr\_abs\_air} to the air and the rest to the cover:
\begin{equation}
Q_{sa}=Q_{st}+s_aQ_{\rm scr},\qquad Q_{\rm cov}=(Q_{\rm roof}-Q_{st})+(1-s_a)Q_{\rm scr},\qquad
Q_{\rm can}+Q_{\rm flr}+Q_{sa}+Q_{\rm cov}+Q_{\rm out}=I_{\rm glob}.
\end{equation}

\subsection{Plant light and photosynthesis}\label{sec:photo}
Stomata and photosynthesis see the PAR on both canopy faces plus lamp PAR:
\begin{equation}
\mathrm{PAR}_p = I_PD_4^P + I_P\rho_2D_5^P + P_{\rm lamp}\,\mathrm{clamp}(U_{12},0,1)\qquad[\text{W\,m}^{-2}],
\end{equation}
($P_{\rm lamp}$ = \code{lamp\_par\_W}). Canopy assimilation (FvCB, constant effective conductance $g=g_sg_m/(g_s+g_m)$) is
\begin{equation}
A = 44010\;\frac{1-e^{-k_pL}}{k_p}\;\min(A_j,A_c,A_p)\quad[\text{mg\,\ce{CO2}\,m}^{-2}\text{s}^{-1}],\qquad k_p=0.7,
\end{equation}
evaluated at $T_2$, $C_1$, $\mathrm{PAR}_p$ with:
\begin{align}
I&=4.6\times10^{-6}\,\mathrm{PAR}_p,\quad C_i=0.509\times10^{-6}C_1,\quad
\Gamma^*=\frac{0.5\,O_a}{S_{c/o,25}\exp\bigl(\frac{T-298.15}{T}\frac{E_{Soc}}{298.15R}\bigr)},\\
J&=\frac{(\alpha I+J_{\max})-\sqrt{(\alpha I+J_{\max})^2-4\theta\alpha IJ_{\max}}}{2\theta},\quad
V_{c\max}=\frac{V_{c\max,25}\,Q_{10,V}^{0.1(T-298.15)}}{1+e^{0.128(T-315.15)}},\\
K_C&=K_{C25}Q_{10,KC}^{0.1(T-298.15)},\quad K_O=K_{O25}Q_{10,KO}^{0.1(T-298.15)},\quad R_d=R_{d,25}\,2^{(T_1-298.15)/10},
\end{align}
and $A_c$, $A_j$ the smaller roots of
\begin{align}
A_c^2-\bigl(V_{c\max}+g(C_i+K_C(1+O_a/K_O))-R_d\bigr)A_c+g\bigl(V_{c\max}(C_i-\Gamma^*)^+-(C_i+K_C(1+O_a/K_O))R_d\bigr)&=0,\\
A_j^2-\bigl(\tfrac14J-R_d+g(C_i+2\Gamma^*)\bigr)A_j+\tfrac14gJ(C_i-\Gamma^*)^+-g(C_i+2\Gamma^*)R_d&=0,
\end{align}
($-R_d$ if the discriminant is negative) and $A_p=1.5V_{c\max}-R_d$.

\subsection{Stomata and transpiration}
With $q_2=P_s(T_1)$ the saturated vapour pressure at the canopy and the day/night switch
$q_7=\bigl(1+e^{(\mathrm{PAR}_p-\delta_1)/\gamma_5}\bigr)^{-1}$ ($\approx1$ by day),
\begin{align}
q_8&=\delta_5+(\delta_4-\delta_5)q_7,\quad q_9=\delta_7+(\delta_6-\delta_7)q_7,\quad
q_{10}=\frac{\mathrm{PAR}_p+\delta_2}{\mathrm{PAR}_p+\delta_3},\\
q_3&=\gamma_4\,q_{10}\,\bigl[1+q_8(C_1T_2/536.4-200)^2\bigr]\bigl[1+q_9(q_2-V_1)^2\bigr]\quad\text{(stomatal resistance, CO}_2\text{ in ppm)},\\
q_1&=\frac{2\rho_3\alpha_5L}{\gamma\gamma_2(\gamma_3+q_3)}\quad(q_2\ge V_1),\qquad
q_1=\frac{2\rho_3\alpha_5L}{\gamma\gamma_2\gamma_3}\quad(q_2<V_1,\ \text{dew}),
\end{align}
transpiration $p_1=q_1(q_2-V_1)$ [kg\,m$^{-2}$\,s$^{-1}$] and its latent heat $l_1=\gamma_2p_1$.

\subsection{Cover node}\label{sec:cover}
The cover has no heat capacity. Linearising its long-wave exchange with the crop air and the sky,
$k_R=4\varepsilon_3\sigma\bigl(\tfrac12(T_2+T_{\rm sky})\bigr)^3$,
\begin{equation}
T_c^{(0)}=\frac{k_R(T_2+T_{\rm sky})+h_{\rm out}T_o+Q_{\rm cov}}{2k_R+h_{\rm out}},\qquad
T_c=T_c^{(0)}+\frac{\gamma_2\,f_{\rm cov}\,6.4\!\times\!10^{-9}\cdot1.7\,|T_2-T_c|^{0.33}\,(V_1-P_s(T_c))^+}{2k_R+h_{\rm out}},
\end{equation}
the second equation solved by four fixed-point passes (latent heat of the vapour condensing on the cover warms it), with
$f_{\rm cov}=1-U_1(1-p_{\rm scr})$ the share of crop air in contact with the cover ($p_{\rm scr}$ = \code{vent\_scr\_perm}).

\subsection{Long-wave radiation}\label{sec:lw}
With $g_1=0.49(1-e^{-\beta_3L})$, $g_2=1-U_1(1-\tau_3)$, $g_3=\tau_1g_2(1-0.49\pi\gamma_1\phi_1)e^{-\beta_3L}$,
$F=\varepsilon_3\sigma(T_2^4-T_c^4)$:
\begin{align}
r_6&=\alpha_3\varepsilon_1\varepsilon_2\,g_1\,\sigma(T_p^4-T_1^4) &&\text{pipe}\to\text{canopy}\\
r_7&=\frac{g_1}{0.49}\varepsilon_2\varepsilon_3\,g_2\,\sigma(T_1^4-T_c^4) &&\text{canopy}\to\text{cover}\\
r_{11}&=\varepsilon_4\,g_3\,\varepsilon_3\sigma(T_5^4-T_c^4) &&\text{floor}\to\text{cover}\\
r_{12}&=\varepsilon_{\rm scr}\,U_1\tau_2\,F &&\text{energy screen (at air }T)\to\text{cover}\\
r_{13}&=A_{\rm cov}\,\varepsilon_6\,F &&\text{air}\to\text{cover}\\
r_{14}&=\varepsilon_4\varepsilon_2(1-0.49\pi\gamma_1\phi_1)\,\sigma(T_1^4-T_5^4) &&\text{canopy}\to\text{floor}
\end{align}

\subsection{Air exchange}\label{sec:vent}
Roof vents: the leeward and windward sides are combined into one equivalent opening,
\begin{align}
B&=\omega_1\nu_6\frac{T_2-T_o}{T_2+T_o},\quad w=n_3v^2,\quad n_3=\nu_2(1-\eta_{11}U_5),\quad
\Phi=\tfrac12\Bigl[U_{\rm lee}^{\,e}\sqrt{(B+c_lw)^+}+g_wU_{\rm win}^{\,e}\sqrt{(B+c_ww)^+}\Bigr],\\
U_8&=\Phi/\sqrt{(B+w)^+}\quad\bigl(\text{if }(B+w)^+>10^{-12};\ \text{else }\tfrac12(U_{\rm lee}^e+g_wU_{\rm win}^e)\bigr),
\end{align}
($e,g_w,c_l,c_w$ = \code{vs\_exp, vs\_gwind, vs\_clee, vs\_cwind}). Air exchange rates [m$^3$\,m$^{-2}$\,s$^{-1}$]:
\begin{align}
f_7&=\frac{U_8\,f_{\rm scr}\,\nu_5n_1}{2\alpha_6}\sqrt{(B+w)^+},\quad f_{\rm scr}=1-U_1(1-p_{\rm scr}),\quad n_1=\nu_1(1-\eta_{10}U_5) &&\text{roof vents}\\
f_4&=\eta_6f_7,\qquad f_2=\eta_6U_1f_5,\quad f_5=\frac{n_1\nu_3U_6\,v\sqrt{n_3}}{2\alpha_6} &&\text{roof, side vents}\\
f_3&=U_7\phi_8/\alpha_6 &&\text{forced ventilation}\\
f_6&=\nu_4\sqrt{\max(v,0.25)^2+c_s|T_2-T_o|} &&\text{leakage (wind + stack)}
\end{align}
The same total $f_2+f_3+f_4+f_6$ carries heat, vapour and \ce{CO2}.

\subsection{Heating}\label{sec:heat}
Where the logged pipe is warm ($T_{\rm log}>T_2+1$) it is used as is; elsewhere a PI controller requests
\begin{equation}\label{eq:pi}
e=U_{11}-T_2,\quad \Delta=K_pe+K_i\max(e_{\rm int},0),\quad T_p=\max\bigl(T_{\rm log},\,T_2+\mathrm{clamp}(\Delta,0,T_{\max}-T_2)\bigr),
\end{equation}
and $\dot e_{\rm int}=e$ except when the logged pipe is warm, the request saturates with $e>0$, or $e_{\rm int}\le0$ with $e<0$ (then 0).
Pipe-to-air convection: $h_4=n_p\,1.99\pi\phi_1\gamma_1|T_p-T_2|^{0.32}(T_p-T_2)$.

\subsection{Lamps}
The full electrical power $Q_l=P_{\rm el}U_{12}$ is released inside: $Q_{lc}=f_{\rm can}(1-e^{-\beta_1L})Q_l$ to the canopy,
$Q_{la}=Q_l-Q_{lc}$ to the air ($P_{\rm el}=117$\,W\,m$^{-2}$, $f_{\rm can}$ = \code{lamp\_fcan}).

\subsection{Sensible and latent heat of the air}
\begin{align}
h_1&=2\alpha_4L(T_1-T_2) &&\text{canopy}\to\text{air}\\
h_2&=\frac{U_2\phi_7}{\alpha_6}\bigl(\rho_3\alpha_5T_o-\gamma_2\rho_3\eta_5(\phi_5-\phi_6)\bigr) &&\text{fan-pad}\\
h_7&=\rho_3\alpha_5(f_2+f_4+f_3+f_6)(T_2-T_o) &&\text{ventilation and leakage}\\
h_{11}&=2\nu_7(T_2-T_5)/(\phi_2+\nu_8) &&\text{air}\to\text{floor conduction}\\
Q_{af}&=k_{\rm ground}(T_2-T_5),\quad Q_{am}=h_m(T_2-T_m) &&\text{air}\to\text{floor, mass convection}\\
l_2&=\gamma_2U_9\phi_9/\alpha_6 &&\text{fog evaporation}
\end{align}

\subsection{Vapour}
\begin{align}
p_2&=\rho_3\frac{U_2\phi_7}{\alpha_6}\bigl(\eta_5(\phi_5-\phi_6)+\phi_6\bigr),\quad p_3=\frac{U_9\phi_9}{\alpha_6},\quad
p_4=\frac{\eta_{12}U_4\lambda_4}{\alpha_6},\quad p_6=\frac{U_2\phi_7}{\alpha_6}\frac{\psi_1}{\omega_2}\frac{V_1}{T_2},\\
p_5&=\frac{\psi_1}{\omega_2}\Bigl(\frac{V_1}{T_2}-\frac{V_o}{T_o}\Bigr)(f_2+f_3+f_4+f_6),\qquad
p_7=U_3\,6.4\!\times\!10^{-9}(V_1-P_s(T_{\rm mc}))^+,\\
p_8&=6.4\!\times\!10^{-9}\cdot1.7\,U_1|T_2-T_{\rm scr}|^{0.33}(V_1-P_s(T_{\rm scr}))^+ &&\text{screen condensation}\\
p_9&=6.4\!\times\!10^{-9}\cdot1200\,|T_1-T_2|^{0.33}(V_1-P_s(T_2))^+ &&\text{supersaturated air}\\
p_{10}&=A_{\rm cov}f_{\rm cov}\,6.4\!\times\!10^{-9}\cdot1.7\,|T_2-T_c|^{0.33}(V_1-P_s(T_c))^+ &&\text{cover condensation}
\end{align}

\subsection{\ce{CO2}}
\begin{equation}
o_1=\frac{\eta_{13}U_4\lambda_4}{\alpha_6},\ o_2=\frac{U_{10}\psi_2}{\alpha_6},\ o_3=\frac{U_2\phi_7}{\alpha_6}(C_1-C_o),\
o_5=(C_1-C_o)(f_2+f_3+f_4+f_6),\ o_R=R_{\rm crop}\,2^{(T_1-293.15)/10}.
\end{equation}

\subsection{Floor, deep soil and internal mass}
$Q_{fd}=k_{\rm deep}(T_5-T_d)$; the floor balance (\ref{eq:T5}) collects air convection and conduction, canopy long-wave,
its own long-wave loss to the cover and its share of the solar column. The internal mass (structure, slabs, water, standing
crop) exchanges $Q_{am}$ with the air and absorbs $Q_{ms}$.

\subsection{Psychrometrics}
$P_s(T)=611.21\exp\bigl[(18.678-\frac{T_c}{234.5})\frac{T_c}{257.14+T_c}\bigr]$ above 0\,$^\circ$C (ice form below),
$T_c=T-273.15$; $\mathrm{RH}=100V_1/P_s(T_2)$.

%======================================================================
\section{Crop model: Aaron V\'elez's generic cucumber model with canopy and LAI}\label{sec:crop}
%======================================================================
The crop model is the generic dry-matter partitioning model of A.~V\'elez (\code{Generic\_Gh\_Model}: \code{models.py},
\code{main.tex}), which casts Marcelis (1994, ``M94'') for cucumber into a generic form, plus his canopy extension
(leaf mass, LAI, retained leaves, pruning). Code: \code{src/crop.jl} (\code{CucumberGeneric}, \code{crop\_step!}) and
\code{src/crop\_measured.jl}. States are per m$^2$ floor; one stem is one modelled plant ($\rho=3.6$ stems\,m$^{-2}$,
Hemming et al.\ 2019). Day $k$ (0 = transplant, 14 August), daily mean air temperature $T$ [$^\circ$C].

\subsection{States}
\begin{center}\small
\begin{tabular}{lll}\toprule
symbol & code & meaning\\\midrule
$W_L$, $W_S$, $W_R$ & \code{W\_leaf, W\_stem, W\_root} & leaf, stem, root dry mass [g\,m$^{-2}$]\\
$L$ & \code{LAI} & leaf area index\\
$N_L$ & \code{leaves} & retained leaves per stem\\
$C_N$ & \code{node} & node counter per stem (fractional progress kept)\\
$(z_i,W_i,n_i)$ & \code{Fruit(dev, W, n, day\_set)} & fruit cohort: stage, dry mass [g\,m$^{-2}$], number [m$^{-2}$]\\
$R$ & \code{reserve} & carbohydrate reserve (DM equivalent) [g\,m$^{-2}$]\\
$Y_{DW}$, $Y_{FW}$ & \code{yield\_DW, yield\_FW} & cumulative harvest [g DM\,m$^{-2}$], [kg FW\,m$^{-2}$]\\
\bottomrule\end{tabular}
\end{center}

\subsection{Generic sub-models (M94 in generic form)}
\begin{align}
a(T)&=\max(\beta_0+\beta_1T,0)\,f_{\rm PAR},\quad \beta_0=-\beta_1\tfrac{0.75}{0.09},\quad
f_{\rm PAR}=\frac{1-e^{-0.5-0.5\,\mathrm{PAR}_k}}{1-e^{-0.5-0.5\cdot6.8}} &&\text{nodes stem}^{-1}\text{d}^{-1}\\
\dot z&=\kappa\max(T-10,0),\qquad \kappa=1/E_{\rm fruit} &&\text{fruit development}\\
P_i&=n_i\,\max(A_0+A_TT,0)\,\frac{u}{(1+u)^p},\quad u=\frac{e^{-s(\min(z_i,1)-z_{\rm peak})}}{p-1} &&\text{fruit sink [g d}^{-1}\text{m}^{-2}]\\
P_V&=\rho\max\bigl(V_{20}+V_T(T-20),0\bigr) &&\text{vegetative demand}
\end{align}
with $s=0.0170\cdot275$, $p=1+1/0.0111$, $z_{\rm peak}=131/275$, $V_T=0.25$, $\mathrm{PAR}_k$ the day's canopy-top PAR
[MJ\,m$^{-2}$\,d$^{-1}$]. The sink amplitude is parametrised by the potential fruit dry mass,
\begin{equation}
A_T=\frac{W_{\rm pot}\,s\,(p-1)}{E_{\rm fruit}},\quad A_0=-10A_T\quad\Rightarrow\quad\int_0^1P_i\frac{dz}{\dot z}\simeq n_iW_{\rm pot},
\end{equation}
so $W_{\rm pot}$ sets the final fruit size and $E_{\rm fruit}$ the duration. A fruit past $z=1$ keeps the sink of $z=1$.

\subsection{Leaf area index}\label{sec:lai}
LAI is a state driven by leaf mass and leaf events; there is no maximum LAI, no recomputation of the standing area and no
senescence; all retained leaves share the canopy's mean mass and area.
\begin{itemize}[nosep]
\item \emph{Leaf events.} Every completed node adds a leaf: $N_L\leftarrow N_L+\lfloor C_{N,k+1}\rfloor-\lfloor C_{N,k}\rfloor$.
\item \emph{Area growth.} $\Delta W_L=f_L\Delta W_V$, $\Delta L=S(d_k)\Delta W_L$ with $f_L=0.59$ and the seasonal SLA
\begin{equation}
S(d)=10^{-4}\Bigl[266+88\sin\frac{2\pi(d+68)}{365}\Bigr]\ \text{m}^2\,\text{g}^{-1},\qquad d_k=226+k,
\end{equation}
(Heuvelink 1999, a tomato regression used provisionally), applied only to new leaf mass.
\item \emph{Pruning.} If $N_L>N_{\rm target}$: $r=N_{\rm target}/N_L$, $W_L\leftarrow rW_L$, $L\leftarrow rL$,
$N_L\leftarrow N_{\rm target}$ (removed mass and area to a ledger; non-leaf mass untouched).
\item \emph{Transplant.} $L_0=0.25$, $C_N=N_0$, $N_L=\lfloor N_0\rfloor$, $W_L=L_0/S(226)$, $W_S+W_R=W_L(1-f_L)/f_L$ split $0.32:0.08$.
\end{itemize}

\subsection{Daily step}
\begin{enumerate}[nosep]
\item \textbf{Supply.} $R_m=Q_{10}^{(T-25)/10}(k_{m,L}W_L+k_{m,S}W_S+k_{m,R}W_R+k_{m,F}\sum_iW_i)$, $B=\max(P_g-R_m,0)/c$, $c=1.4$.
\item \textbf{Allocation} with reserve: $\Delta W_j=\frac{P_j}{P}\min(B+R,P)$, $P=P_V+\sum_iP_i$, $R\leftarrow B+R-\sum_j\Delta W_j$.
\item \textbf{Vegetative growth and LAI}: as above; non-leaf mass split stem:root $=0.32:0.08$; $W_i\mathrel{+}=\Delta W_i$.
\item \textbf{Development and appearance} (before topping, day 104): $z_i\mathrel{+}=\dot z$, $C_N\mathrel{+}=a(T)$; nodes $\ge5$ carry a
flower; from day 15 the flowers set $n_{\rm new}=n_{\rm fl}\,\rho\,\phi_k\,\mathrm{clamp}\bigl(\frac{B/P-r_{\rm low}}{r_{\rm high}-r_{\rm low}},0,1\bigr)$
fruits (new cohort at $z=0$, growing from the next day), $\phi_k=1-\text{Pruning}_k/100$ the grower's kept share.
\item \textbf{Harvest} when $W_i/n_i\ge323\,\mathrm{DMC}_{\rm fruit}$ (observed mean A+B fruit, 323\,g FW) or $z_i\ge3$.
\item \textbf{Pruning} to $N_{\rm target}$.
\end{enumerate}

\subsection{Light and photosynthesis of the crop path}
Gross assimilation $P_g$ [g\,CH$_2$O\,m$^{-2}$\,d$^{-1}$] is the FvCB model of Section~\ref{sec:photo} summed over the 5-min
samples of the day, at measured air T and \ce{CO2}, with the organisers' canopy-top PAR
\begin{equation}
\mathrm{PAR}=\bigl[\mathrm{PAR}_{\rm out}\cdot0.65\,(1-0.98U_{\rm bl})(1-0.25U_1)+187\,U_{12}\bigr]/4.6\quad[\text{W\,m}^{-2}],
\end{equation}
multiplied by the both-faces factor $g(L)=(D_4^P+\rho_2D_5^P)/t_r$ of the PAR column with screens retracted, effective leaf
area $(1-e^{-0.75L})/0.75$ and constant $R_d$.

\subsection{Climate inputs from the crop}
\code{crop\_series}: $L$; $R_{\rm crop}=\bigl[2^{-0.5}(k_{m,S}W_S+k_{m,R}W_R+k_{m,F}\sum W_i)+(c-1)\Delta DM\bigr]\frac{44}{30}\frac{1000}{86400}$
(non-leaf maintenance at 20\,$^\circ$C + growth respiration; leaf dark respiration is in the FvCB model);
$\mathrm{FW}=(W_S/0.06+W_R/0.07+\sum W_i/\mathrm{DMC}_{\rm fruit})/1000$.


%======================================================================
\section{Parameters}\label{sec:params}
%======================================================================
{\small
\begin{longtable}{>{\raggedright\arraybackslash}p{0.17\textwidth}>{\raggedright\arraybackslash}p{0.17\textwidth}>{\raggedright\arraybackslash}p{0.13\textwidth}>{\raggedright\arraybackslash}p{0.43\textwidth}}
\caption{Climate and photosynthesis parameters (\texttt{configfiles/climate\_n79.json}). Inferred parameters: [min, centre, max] of the uniform prior; the centre is the default value.}\label{tab:pclim}\\
\toprule
name & value (prior) & units & meaning; source\\
\midrule\endfirsthead
\toprule name & value (prior) & units & meaning; source\\ \midrule\endhead
\multicolumn{4}{l}{\textit{geometry}}\\
\texttt{alpha3} & 0.3 & m$^{2}$ m$^{-2}$ & pipe surface seen by the canopy (long-wave); V2.0 (audit: geometric 0.16)\\
\texttt{alpha6} & 10000 & m$^{2}$ & floor area used to express absolute capacities per m$^{2}$ floor; V2.0\\
\texttt{cover\_area} & 1.6 & m$^{2}$ m$^{-2}$ & exterior cover area per floor area (roof + one side wall); Note 67\\
\texttt{gamma1} & 1 & m m$^{-2}$ & heating-pipe length per floor area (rail pipe, one circuit); Note 78\\
\texttt{n\_pipes} & 1 & - & number of heating circuits (crop pipe folded into the rail pipe by the inputs); Note 54/79\\
\texttt{nu3} & 2352.9 & m$^{2}$ & side-wall vent surface (side vents, input U6); V2.0\\
\texttt{nu5} & 3000 & m$^{2}$ & roof-vent area (0.30 of the floor, insect nets); Notes 78, 92\\
\texttt{nu6} & 1 & m & vertical dimension of a roof-vent opening; V2.0\\
\texttt{nu8} & 0.64 & m & floor-to-ground distance in the air-floor conduction length; V2.0\\
\texttt{phi1} & 0.051 & m & outer diameter of the heating pipe; V2.0 / Note 78\\
\texttt{phi2} & 5.4 & m & mean air height (volume per floor area); Note 67 (compartment drawing)\\
\multicolumn{4}{l}{\textit{constants}}\\
\texttt{alpha5} & 1005 & J kg$^{-1}$ K$^{-1}$ & specific heat of air; Note 78\\
\texttt{gamma} & 65.8 & Pa K$^{-1}$ & psychrometric constant; V2.0\\
\texttt{gamma2} & $2.45\times10^{6}$ & J kg$^{-1}$ & latent heat of evaporation; V2.0\\
\texttt{omega1} & 9.81 & m s$^{-2}$ & gravity; V2.0\\
\texttt{omega2} & 8314 & J kmol$^{-1}$ K$^{-1}$ & molar gas constant; V2.0\\
\texttt{psi1} & 18 & kg kmol$^{-1}$ & molar mass of water; V2.0\\
\texttt{rho3} & 1.2 & kg m$^{-3}$ & air density; V2.0\\
\texttt{sigma} & $5.67\times10^{-8}$ & W m$^{-2}$ K$^{-4}$ & Stefan-Boltzmann constant; V2.0\\
\multicolumn{4}{l}{\textit{air}}\\
\texttt{air\_cap\_mult} & \textbf{inferred} [1, 3, 12] & - & effective air heat-capacity multiplier (fast masses lumped with the air); n79 prior (Note 78; f\_mass Note 104)\\
\multicolumn{4}{l}{\textit{canopy}}\\
\texttt{alpha1} & 1500 & J K$^{-1}$ m$^{-2}$ leaf & canopy heat capacity per leaf area; V2.0\\
\texttt{alpha4} & 5 & W m$^{-2}$ K$^{-1}$ & leaf-air convective coefficient (per leaf side); Vanthoor 2011\\
\texttt{beta1} & 0.7 & - & canopy PAR extinction coefficient; V2.0\\
\texttt{beta3} & 0.27 & - & canopy NIR extinction coefficient (also used for the long-wave terms); V2.0 (audit: K\_FIR 0.94)\\
\texttt{delta1} & 5 & W m$^{-2}$ & PAR defining day/night for the stomatal resistance; V2.0\\
\texttt{delta2} & 4.3 & W m$^{-2}$ & stomatal light response parameter; V2.0\\
\texttt{delta3} & 0.54 & W m$^{-2}$ & stomatal light response parameter; V2.0\\
\texttt{delta4} & $6.1\times10^{-7}$ & ppm$^{-2}$ & stomatal CO2 response, day; V2.0\\
\texttt{delta5} & $1.1\times10^{-11}$ & ppm$^{-2}$ & stomatal CO2 response, night; V2.0\\
\texttt{delta6} & $4.3\times10^{-6}$ & Pa$^{-2}$ & stomatal vapour-pressure response, day; V2.0\\
\texttt{delta7} & $5.2\times10^{-6}$ & Pa$^{-2}$ & stomatal vapour-pressure response, night; V2.0\\
\texttt{epsil2} & 1 & - & canopy long-wave emissivity; V2.0\\
\texttt{gamma3} & 275 & s m$^{-1}$ & canopy boundary-layer resistance to vapour; V2.0\\
\texttt{gamma4} & \textbf{inferred} [40, 82, 200] & s m$^{-1}$ & minimum stomatal resistance; n79 prior (Note 78; f\_mass Note 104)\\
\texttt{gamma5} & -1 & W m$^{-2}$ & inverse slope of the day/night switch of the stomatal resistance; V2.0\\
\texttt{k\_lai\_photo} & 0.7 & - & extinction in the effective leaf area (1-exp(-kL))/k of the climate photosynthesis; V2.2 K\_EXT\\
\texttt{rho1} & 0.07 & - & canopy PAR reflection at high LAI; V2.0\\
\texttt{rho\_can\_nir} & 0.35 & - & canopy NIR reflection at high LAI; Vanthoor 2011 / Note 103\\
\multicolumn{4}{l}{\textit{photosynthesis}}\\
\texttt{E\_Soc} & -24460 & J mol$^{-1}$ & activation energy of Sco; V2.0\\
\texttt{J\_max} & $2\times10^{-4}$ & mol m$^{-2}$ s$^{-1}$ & maximum electron transport rate (climate model; the crop calibration infers its own); V2.0\\
\texttt{K\_C25} & $3\times10^{-4}$ & mol m$^{-3}$ & Michaelis constant for CO2 at 25 C; V2.0\\
\texttt{K\_O25} & 0.3 & mol m$^{-3}$ & Michaelis constant for O2 at 25 C; V2.0\\
\texttt{O\_a} & 0.21 & mol m$^{-3}$ & O2 concentration; V2.0\\
\texttt{Q10\_KC} & 2.1 & - & Q10 of K\_C; V2.0\\
\texttt{Q10\_KO} & 1.2 & - & Q10 of K\_O; V2.0\\
\texttt{Q10\_Vcmax} & 2.4 & - & Q10 of Vcmax; V2.0\\
\texttt{Rd\_day} & $1\times10^{-6}$ & mol m$^{-2}$ s$^{-1}$ & leaf dark respiration at 25 C (x Q10 2 on canopy T); V2.0 / Note 86\\
\texttt{Rgas} & 8.31447 & J mol$^{-1}$ K$^{-1}$ & gas constant; V2.0\\
\texttt{Sco25} & 2800 & - & Rubisco CO2/O2 specificity at 25 C; V2.0\\
\texttt{V\_cmax25} & $1.5\times10^{-4}$ & mol m$^{-2}$ s$^{-1}$ & maximum carboxylation rate at 25 C; V2.0\\
\texttt{alpha} & 0.35 & - & quantum efficiency of electron transport; V2.0\\
\texttt{g\_m} & 0.25 & mol m$^{-2}$ s$^{-1}$ & mesophyll conductance in the constant effective conductance; Bush 2023 (V2.0)\\
\texttt{g\_s} & 0.3 & mol m$^{-2}$ s$^{-1}$ & stomatal conductance in the constant effective conductance; Bush 2023 (V2.0)\\
\texttt{theta} & 0.7 & - & curvature of the light response; V2.0\\
\multicolumn{4}{l}{\textit{floor}}\\
\texttt{C\_soil} & \textbf{inferred} [$1.5\times10^{5}$, $6\times10^{5}$, $2\times10^{6}$] & J m$^{-2}$ K$^{-1}$ & heat capacity of the active floor layer; n79 prior (Note 78; f\_mass Note 104)\\
\texttt{alpha7} & 0.5 & - & floor NIR absorption (reflection 1 - alpha7); V2.0\\
\texttt{epsil4} & 1 & - & floor long-wave emissivity; V2.0\\
\texttt{k\_deep} & \textbf{inferred} [0.5, 1.5, 3] & W m$^{-2}$ K$^{-1}$ & floor to deep-soil conductance; n79 prior (Note 78; f\_mass Note 104)\\
\texttt{k\_ground} & \textbf{inferred} [1, 8, 40] & W m$^{-2}$ K$^{-1}$ & air-floor convective conductance; n79 prior (Note 78; f\_mass Note 104)\\
\texttt{nu7} & 0.85 & W m$^{-1}$ K$^{-1}$ & soil conductivity in the air-floor conduction h11; V2.0\\
\texttt{rho2} & 0.65 & - & floor PAR reflection; V2.0\\
\texttt{sd\_tau} & $1.2096\times10^{6}$ & s & time constant of the deep-soil temperature; Note 47\\
\multicolumn{4}{l}{\textit{mass}}\\
\texttt{C\_mass} & \textbf{inferred} [20000, 50000, $1.2\times10^{5}$] & J m$^{-2}$ K$^{-1}$ & heat capacity of the internal mass (structure, slabs, water); n79 prior (Note 78; f\_mass Note 104)\\
\texttt{crop\_cp} & 3800 & J kg$^{-1}$ K$^{-1}$ & specific heat of fresh plant tissue; Note 91\\
\texttt{f\_mass} & \textbf{inferred} [0, 0.3, 1] & - & share of the floor-level solar absorbed by the internal mass; n79 prior (Note 78; f\_mass Note 104)\\
\texttt{h\_mass} & \textbf{inferred} [1, 5, 40] & W m$^{-2}$ K$^{-1}$ & air-internal mass heat-transfer coefficient; n79 prior (Note 78; f\_mass Note 104)\\
\multicolumn{4}{l}{\textit{cover}}\\
\texttt{cov\_abs} & 0.1 & - & share of the global radiation absorbed by the glass (booked to the cover); Note 46\\
\texttt{epsil3} & 1 & - & sky/cover long-wave emissivity; V2.0\\
\texttt{epsil6} & 0.44 & - & cover long-wave emissivity (air-cover term r13); V2.0\\
\texttt{eta1} & 0.1 & - & share of the global radiation absorbed by the structure (booked to the air); V2.0\\
\texttt{eta2} & 0.5 & - & PAR share of the global radiation; V2.0\\
\texttt{eta3} & 0.5 & - & NIR share of the global radiation; V2.0\\
\texttt{h\_cover\_out} & \textbf{inferred} [10, 25, 50] & W m$^{-2}$ K$^{-1}$ & cover-outside convective coefficient; n79 prior (Note 78; f\_mass Note 104)\\
\texttt{tau1} & 0.722222 & - & cover PAR transmission; (1-eta1) tau1 = 0.65 roof transmission; Note 78 (organisers 0.60-0.70)\\
\texttt{tau2} & 1 & - & cover long-wave transmission in the screen term r12; V2.0\\
\multicolumn{4}{l}{\textit{screens}}\\
\texttt{eps\_screen} & 0.4 & - & energy-screen long-wave emissivity (upper face, r12); V2.2\\
\texttt{scr\_abs\_air} & 0.5 & - & share of screen-absorbed solar released to the crop air (rest to the cover); Note 101 (PROVISIONAL)\\
\texttt{scr\_bl\_rho\_nir} & 0.5 & - & blackout screen NIR reflection (PROVISIONAL, = PAR); Note 101\\
\texttt{scr\_bl\_rho\_par} & 0.5 & - & blackout screen PAR reflection (PROVISIONAL, assumed); Note 101\\
\texttt{scr\_bl\_tau\_nir} & 0.02 & - & blackout screen NIR transmission (PROVISIONAL, = PAR); Note 101\\
\texttt{scr\_bl\_tau\_par} & 0.02 & - & blackout screen PAR transmission at full closure (PROVISIONAL); organisers 98 \% interception\\
\texttt{scr\_es\_rho\_nir} & 0.15 & - & energy screen NIR reflection (PROVISIONAL, = PAR); Note 101\\
\texttt{scr\_es\_rho\_par} & 0.15 & - & energy screen PAR reflection (PROVISIONAL, assumed); Note 101\\
\texttt{scr\_es\_tau\_nir} & 0.75 & - & energy screen NIR transmission (PROVISIONAL, = PAR); Note 101\\
\texttt{scr\_es\_tau\_par} & 0.75 & - & energy screen PAR transmission at full closure (PROVISIONAL); organisers 25 \% interception\\
\texttt{tau3} & 0.11 & - & energy-screen long-wave transmission; V2.0\\
\multicolumn{4}{l}{\textit{heating}}\\
\texttt{epsil1} & 0.88 & - & pipe long-wave emissivity; V2.0\\
\texttt{hc\_Ki} & 0.00555556 & s$^{-1}$ & PI heating controller integral gain (tau\_i 30 min); Note 34\\
\texttt{hc\_Kp} & 10 & K K$^{-1}$ & PI heating controller proportional gain; Note 34\\
\texttt{hc\_Tmax} & 338.15 & K & maximum pipe temperature (65 C); Note 33\\
\multicolumn{4}{l}{\textit{ventilation}}\\
\texttt{eta10} & 0 & - & shading effect on the discharge coefficient; V2.0\\
\texttt{eta11} & 0 & - & shading effect on the wind-pressure coefficient; V2.0\\
\texttt{eta6} & 0.551 & - & ventilation reduction factor (Vanthoor 8.70, zeta 0.33); V2.0\\
\texttt{leak\_cs} & \textbf{inferred} [0, 0.3, 1] & m$^{2}$ s$^{-2}$ K$^{-1}$ & stack (temperature-difference) term of the leakage; n79 prior (Note 78; f\_mass Note 104)\\
\texttt{nu1} & \textbf{inferred} [0.3, 0.65, 1] & - & roof-vent discharge coefficient; n79 prior (Note 78; f\_mass Note 104)\\
\texttt{nu2} & \textbf{inferred} [0.005, 0.05, 0.3] & - & roof-vent wind-pressure coefficient; n79 prior (Note 78; f\_mass Note 104)\\
\texttt{nu4} & \textbf{inferred} [$1\times10^{-5}$, $4.5\times10^{-5}$, $1\times10^{-4}$] & m s$^{-1}$ per m s$^{-1}$ & leakage coefficient; n79 prior (Note 78; f\_mass Note 104)\\
\texttt{vent\_scr\_perm} & \textbf{inferred} [0, 0.4, 1] & - & share of roof ventilation and cover contact reaching the crop air through a closed energy screen; n79 prior (Note 78; f\_mass Note 104)\\
\texttt{vs\_clee} & 1 & - & leeward wind-pressure multiplier; Note 58\\
\texttt{vs\_cwind} & \textbf{inferred} [0.5, 1, 10] & - & windward wind-pressure multiplier; n79 prior (Note 78; f\_mass Note 104)\\
\texttt{vs\_exp} & \textbf{inferred} [0.5, 0.9, 1.2] & - & curvature of opening -> flow; n79 prior (Note 78; f\_mass Note 104)\\
\texttt{vs\_gwind} & \textbf{inferred} [0.5, 1.5, 4] & - & windward discharge multiplier relative to leeward; n79 prior (Note 78; f\_mass Note 104)\\
\multicolumn{4}{l}{\textit{lamps}}\\
\texttt{lamp\_Pel} & 117 & W m$^{-2}$ & lamp electrical power at 100 \% (187 umol / 1.6 umol J$^{-1}$); Note 90\\
\texttt{lamp\_fcan} & \textbf{inferred} [0.05, 0.3, 0.9] & - & share of lamp power absorbed by a closed canopy; n79 prior (Note 78; f\_mass Note 104)\\
\texttt{lamp\_par\_W} & 40.65 & W m$^{-2}$ & lamp PAR at the canopy at 100 \% (187 umol / 4.6); Note 71\\
\multicolumn{4}{l}{\textit{actuators}}\\
\texttt{eta12} & $4.43\times10^{-8}$ & kg J$^{-1}$ & vapour released per J of direct-heater heat (U4); V2.0\\
\texttt{eta13} & 0.057 & mg J$^{-1}$ & CO2 released per J of direct-heater heat (U4); V2.0\\
\texttt{eta5} & 0.5 & - & fan-pad efficiency (U2); V2.0\\
\texttt{lamb4} & $5\times10^{5}$ & W & direct air heater capacity (U4); V2.0\\
\texttt{phi5} & 0.014 & kg kg$^{-1}$ & vapour content of fan-pad air (U2); V2.0\\
\texttt{phi6} & 0.0079 & kg kg$^{-1}$ & vapour content of outside air for the fan-pad (U2); V2.0\\
\texttt{phi7} & 16.7 & m$^{3}$ s$^{-1}$ & fan-pad air flow capacity (U2); V2.0\\
\texttt{phi8} & 666.6 & m$^{3}$ s$^{-1}$ & forced-ventilation capacity (U7); V2.0\\
\texttt{phi9} & 0.916 & kg s$^{-1}$ & fog capacity (U9); V2.0\\
\texttt{psi2} & 35556 & mg s$^{-1}$ & CO2 valve capacity (12.8 g m$^{-2}$ h$^{-1}$ x alpha6); Note 78\\
\multicolumn{4}{l}{\textit{noise}}\\
\texttt{sigma\_C} & \textbf{inferred} [10, 60, 400] & mg m$^{-3}$ & likelihood noise; n79 prior\\
\texttt{sigma\_RH} & \textbf{inferred} [0.5, 6, 25] & \% & likelihood noise; n79 prior\\
\texttt{sigma\_T} & \textbf{inferred} [0.2, 1.5, 6] & K & likelihood noise; n79 prior\\
\bottomrule
\end{longtable}}

{\small
\begin{longtable}{>{\raggedright\arraybackslash}p{0.17\textwidth}>{\raggedright\arraybackslash}p{0.17\textwidth}>{\raggedright\arraybackslash}p{0.13\textwidth}>{\raggedright\arraybackslash}p{0.43\textwidth}}
\caption{Crop parameters (\texttt{configfiles/crop\_n105.json}).}\label{tab:pcrop}\\
\toprule
name & value (prior) & units & meaning; source\\
\midrule\endfirsthead
\toprule name & value (prior) & units & meaning; source\\ \midrule\endhead
\multicolumn{4}{l}{\textit{M94}}\\
\texttt{E\_fruit} & \textbf{inferred} [150, 275, 450] & C d & fruit duration (degree-days above 10 C); n94/n105 prior (Notes 74, 94, 98)\\
\texttt{V20} & \textbf{inferred} [1, 3.2, 8] & g DM stem$^{-1}$ d$^{-1}$ & vegetative demand at 20 C; n94/n105 prior (Notes 74, 94, 98)\\
\texttt{VT} & 0.25 & g DM stem$^{-1}$ d$^{-1}$ C$^{-1}$ & temperature slope of vegetative demand; Marcelis 1994\\
\texttt{W\_pot} & \textbf{inferred} [20, 60.7, 120] & g DM fruit$^{-1}$ & potential fruit dry mass; n94/n105 prior (Notes 74, 94, 98)\\
\texttt{app\_base} & 8.33333 & C & base temperature of node appearance (beta0 = -beta1 x 0.75/0.09); Marcelis 1994\\
\texttt{app\_beta1} & \textbf{inferred} [0.05, 0.0882, 0.13] & nodes stem$^{-1}$ d$^{-1}$ C$^{-1}$ & node (flower) appearance slope; n94/n105 prior (Notes 74, 94, 98)\\
\texttt{dev\_base} & 10 & C & base temperature of fruit development and of the sink amplitude; Marcelis 1994\\
\texttt{par\_ref} & 6.8 & MJ m$^{-2}$ d$^{-1}$ & PAR at which M94 fitted the appearance rate; Marcelis 1994\\
\texttt{sink\_p} & 91.0901 & - & fruit sink exponent p = 1 + 1/0.0111; Marcelis 1994 via Aaron\\
\texttt{sink\_peak} & 0.476364 & - & stage of peak sink = 131/275; Marcelis 1994 via Aaron\\
\texttt{sink\_s} & 4.675 & - & fruit sink shape s = 0.0170 x 275; Marcelis 1994 via Aaron\\
\multicolumn{4}{l}{\textit{canopy}}\\
\texttt{LAI0} & 0.25 & m$^{2}$ m$^{-2}$ & LAI at transplant (assumed); Note 94\\
\texttt{frac\_root} & 0.08 & - & root weight in the stem:root split; V2.2\\
\texttt{frac\_stem} & 0.32 & - & stem weight in the stem:root split; V2.2\\
\texttt{gen\_N0} & \textbf{inferred} [2, 4, 8] & nodes & nodes at transplant; n94/n105 prior (Notes 74, 94, 98)\\
\texttt{gen\_doy0} & 226 & day of year & transplant day (14 Aug); Note 94\\
\texttt{gen\_fL} & 0.59 & - & leaf share of vegetative growth; Aaron (Note 94)\\
\texttt{gen\_leaf\_target} & \textbf{inferred} [20, 25, 30] & leaves stem$^{-1}$ & retained-leaf target (pruning); n94/n105 prior (Notes 74, 94, 98)\\
\texttt{sla\_amp} & 88 & cm2 g$^{-1}$ & seasonal SLA amplitude; Heuvelink 1999 eq. 3\\
\texttt{sla\_mean} & 266 & cm2 g$^{-1}$ & seasonal SLA mean; Heuvelink 1999 eq. 3\\
\texttt{sla\_phase} & 68 & d & seasonal SLA phase; Heuvelink 1999 eq. 3\\
\texttt{stem\_density} & 3.6 & stems m$^{-2}$ & stem density (one stem = one modelled plant); Hemming et al. 2019 (Note 92)\\
\multicolumn{4}{l}{\textit{architecture}}\\
\texttt{gen\_flower\_node} & 5 & node & first node with a flower; Aaron (Note 94)\\
\texttt{set\_start\_day} & 15 & d & first day with fruit set; Note 66\\
\texttt{top\_day} & 104 & d & topping: no node appearance after this day; V2.2\\
\multicolumn{4}{l}{\textit{respiration}}\\
\texttt{Q10\_resp} & 2 & - & Q10 of maintenance; V2.2\\
\texttt{asrq} & 1.4 & g CH2O g$^{-1}$ DM & construction cost; Aaron (1.4)\\
\texttt{k\_m\_fruit} & \textbf{inferred} [0.003, 0.01, 0.02] & g g$^{-1}$ d$^{-1}$ & fruit maintenance at 25 C; n94/n105 prior (Notes 74, 94, 98)\\
\texttt{k\_m\_leaf} & \textbf{inferred} [0.01, 0.03, 0.06] & g g$^{-1}$ d$^{-1}$ & leaf maintenance at 25 C; n94/n105 prior (Notes 74, 94, 98)\\
\texttt{k\_m\_root} & 0.01 & g g$^{-1}$ d$^{-1}$ & root maintenance at 25 C; V2.2\\
\texttt{k\_m\_stem} & \textbf{inferred} [0.005, 0.015, 0.03] & g g$^{-1}$ d$^{-1}$ & stem maintenance at 25 C; n94/n105 prior (Notes 74, 94, 98)\\
\multicolumn{4}{l}{\textit{fruit set}}\\
\texttt{set\_r\_high} & \textbf{inferred} [0.2, 1, 2.5] & - & source/sink ratio with full fruit set; n94/n105 prior (Notes 74, 94, 98)\\
\texttt{set\_r\_low} & \textbf{inferred} [0, 0.3, 1] & - & source/sink ratio with no fruit set; n94/n105 prior (Notes 74, 94, 98)\\
\multicolumn{4}{l}{\textit{quality}}\\
\texttt{DMC\_fruit} & \textbf{inferred} [0.025, 0.035, 0.045] & - & fruit dry-matter content; n94/n105 prior (Notes 74, 94, 98)\\
\multicolumn{4}{l}{\textit{harvest}}\\
\texttt{harvest\_FW} & 323 & g FW & picking size (observed mean A+B fruit); Note 92\\
\multicolumn{4}{l}{\textit{photosynthesis}}\\
\texttt{J\_max} & \textbf{inferred} [$1\times10^{-4}$, $2\times10^{-4}$, $3.5\times10^{-4}$] & mol m$^{-2}$ s$^{-1}$ & maximum electron transport rate (crop path); n94/n105 prior (Notes 74, 94, 98)\\
\texttt{k\_ext} & 0.75 & - & extinction in the effective leaf area of the crop path; V2.2 (climate uses 0.7)\\
\multicolumn{4}{l}{\textit{light}}\\
\texttt{lamp\_umol} & 187 & umol m$^{-2}$ s$^{-1}$ & lamp PAR at 100 \%; AGC ReadMe\\
\texttt{transmission} & 0.65 & - & greenhouse PAR transmission in the measured-light formula; Note 78\\
\multicolumn{4}{l}{\textit{series}}\\
\texttt{dmc\_root} & 0.07 & - & root dry-matter content (standing fresh mass); Note 91 (assumed)\\
\texttt{dmc\_stem} & 0.06 & - & stem dry-matter content (standing fresh mass); Note 91 (assumed)\\
\multicolumn{4}{l}{\textit{noise}}\\
\texttt{sigma\_fd} & \textbf{inferred} [0.5, 3, 15] & - & likelihood noise; n94/n105 prior (Notes 74, 94, 98)\\
\texttt{sigma\_leaf} & \textbf{inferred} [0.5, 3, 20] & - & likelihood noise; n94/n105 prior (Notes 74, 94, 98)\\
\texttt{sigma\_n} & \textbf{inferred} [0.3, 3, 20] & - & likelihood noise; n94/n105 prior (Notes 74, 94, 98)\\
\texttt{sigma\_y} & \textbf{inferred} [0.02, 0.3, 3] & - & likelihood noise; n94/n105 prior (Notes 74, 94, 98)\\
\bottomrule
\end{longtable}}


%======================================================================
\section{Parameter inference}\label{sec:inference}
%======================================================================
Both calibrations use the t-walk sampler, uniform priors (Tables~\ref{tab:pclim}, \ref{tab:pcrop}), a Gaussian likelihood
with inferred noise per observable,
\begin{equation}
U(\theta,\sigma)=\sum_k\Bigl[n_k\log\sigma_k+\frac{\mathrm{SSE}_k(\theta)}{2\sigma_k^2}\Bigr],
\end{equation}
two or more seeds, split-$\hat R$ and an explicit burn-in.
\begin{itemize}[nosep]
\item \textbf{Crop} (\code{test/calibrate\_crop.jl}): 13 parameters + 4 noise terms; observations weekly yield, weekly
harvested fruit number, leaves per stem, fruit development time; driver measured climate.
\item \textbf{Climate} (\code{test/calibrate\_climate.jl}): 18 parameters + 3 noise terms; 7-day blocks with a spin-up day and the
measured initial state; TRAIN days 7--13, 28--34, 63--69, 91--97, 110--114; VAL 42--48, 77--83; TEST the rest; observations
hourly means of air T and RH and dark-hour \ce{CO2}; vents, pipe, screens, lamps and \ce{CO2} valve are logged inputs.
\end{itemize}
\paragraph{What to infer and what to fix.} Measurable or documented quantities are fixed (geometry, materials, actuator
capacities, crop counts and dates); missing values of this kind are questions for the organisers or data sheets. Only
coefficients of approximate sub-models are inferred. No bias or measurement-offset terms. As the physics becomes more
descriptive, effective parameters should be retired: $h_{\rm out}\to2.8+1.2v$ with a conserving cover node;
\code{air\_cap\_mult} $\to$ explicit fast masses; \code{lamp\_fcan} $\to$ lamp spectrum through the column;
\code{vent\_scr\_perm} $\to$ screen permeability in a top compartment.


%======================================================================
\section{Known structural gaps}\label{sec:gaps}
%======================================================================
These are properties of the equations above (read-only audit of 21 October 2026 and the project notes); they are the
starting list for the physics work.
\begin{enumerate}[nosep]
\item \textbf{Cover node not conserving.} The node of Section~\ref{sec:cover} balances $k_R(T_2-T_c)$ only; it does not
receive $r_7$, $r_{11}$, $r_{12}$, $r_{13}$ nor pipe radiation (40\,W\,m$^{-2}$ residual on a test night). $r_{11}$ contains the PAR
transmission $\tau_1$. The node is not scaled by the cover area while $r_{13}$ is.
\item \textbf{No top compartment.} The screen sits at $\tfrac12(T_o+T_2)$ for condensation and at $T_2$ in $r_{12}$; $r_{13}$ has
no screen factor; cover condensation draws on the crop air (through $f_{\rm cov}$). Geometry: screen $\approx4.5$\,m, ridge
$\approx5.8$\,m, $\approx0.9$\,m$^3$\,m$^{-2}$ above the screen.
\item \textbf{Long-wave network against LAI.} $r_{14}$ has no canopy fraction; $\beta_3$ (NIR) is used for long-wave instead of
$K_{\rm FIR}=0.94$; pipes radiate only to the canopy ($\alpha_3=0.3$ vs geometric 0.16), no pipe--floor, pipe--cover terms.
\item \textbf{Lamps.} At $f_{\rm can}=0.3$, $L=2.3$ the canopy heat (28.1\,W\,m$^{-2}$) is below the absorbed lamp PAR (34.3); no
lamp light reaches the floor.
\item \textbf{Inherited V2.0 forms}: side-vent flow scaled by the energy screen ($f_2=\eta_6U_1f_5$); air--floor conduction length
uses the air height ($\phi_2+\nu_8$); mechanical cooling removes vapour but no heat; $h_2$ fan-pad form. Zero input in AiCU
except $h_{11}$.
\item \textbf{Fit at the prior centre} (V2.2 smoke 4): air $\approx1.2$\,K too cold (day $-1.6$, night $-1.0$), RH $\approx-10\,\%$.
\item \textbf{Photosynthesis}: big leaf; $J_{\max}$ of the climate model (JSON) differs from the crop posterior; $k=0.7$ (climate) vs
0.75 (crop path); constant $R_d$ in the crop path.
\item \textbf{Crop}: no fruit abortion (fruit number 107 vs 116 observed); SLA from tomato; bimodal fruit size/duration posterior.
\item \textbf{Provisional optics}: screen reflectances and long-wave properties, NIR = PAR for the screens, $\rho_\infty^{\rm NIR}$,
\code{scr\_abs\_air}.
\item \textbf{Coupling}: crop $\to$ climate only, offline.
\item \textbf{Numerical noise of the likelihood} (Section~\ref{sec:verify}): at reltol $10^{-4}$ the energy changes by up to $\sim$60
units for a $10^{-14}$ relative change of one parameter. Smoothing the switches (heating, clamps, condensation onsets),
stopping the solver at input steps, or a tighter tolerance are needed before long MCMC runs.
\end{enumerate}


%======================================================================
\section{Verification}\label{sec:verify}
%======================================================================
The code was obtained from GreenhouseSim V2.2 (git tag \code{v2.2-flags-final}) by resolving every flag to its n79 value
(\code{PRIORS=n79}, \code{SOLARC=1}, \code{PIPE2=1}, \code{CO2NIGHT=1}) and the crop to the n105 set-up, and deleting the
branches not taken. Four checks support the claim that the code is n79 and that this document is the code:
\begin{enumerate}[nosep]
\item \code{tools/audit.py}: no flag name or \code{hasproperty(p, ...)} switch is left in \code{src/}; every parameter the code
reads is in the JSON files and every JSON key is read.
\item \code{test/equivalence\_synthetic.jl} (no data): V2.2 with the n79 flags and GreenhouseTwin on the same synthetic
inputs: \code{rhs!} at 5000 random states and times (relative difference $<10^{-10}$), 3-day trajectories, solar fluxes, and the
crop over 116 synthetic days.
\item \code{test/equivalence\_n79.jl} (AGC data): both codes build their inputs independently from the AGC files (compared
array by array), parameters, \code{rhs!} at random states, TRAIN blocks, likelihood energies at the prior centre and at a
jittered point, the full season, the solar budget, the crop at the n105 prior centre, timing and allocations.
\item Term dump: \code{test/dump\_terms.jl} writes every term of \code{model\_terms} at test states; \code{tools/check\_tex.py}
evaluates the equations of Section~\ref{sec:climate} written in Python from this document and compares.
\end{enumerate}
\paragraph{Results (21 October 2026, Julia 1.10).} Audit: pass. Synthetic: 35/35. Document against code (\code{check\_tex.py},
400 states, 60 terms): largest relative difference $2\times10^{-14}$; this check found and corrected a sign error in the FvCB
quadratics of an earlier draft of Section~\ref{sec:photo}. AGC data: all 20 input arrays and 112 shared parameters identical;
\code{rhs!} derivatives within $3\times10^{-16}$ relative; solar fluxes and crop series identical (crop to $10^{-14}$);
\code{rhs!} 1.7\,$\mu$s against 2.3\,$\mu$s for V2.2, no allocations.

\paragraph{Numerical noise of the trajectories and of the likelihood.} Because the right-hand side is discontinuous in
time (5-min step inputs) and in state (heating switch $T_{\rm log}>T_2+1$, PI clamp, $\max$/$\min$, condensation onsets),
the adaptive solver turns rounding-level differences into finite differences of the solution. V2.2 compared with itself
after scaling one parameter by $1+10^{-14}$:
\begin{center}\small
\begin{tabular}{lccc}\toprule
reltol & TRAIN-block $T_2$ difference & energy sd (10 perturbations) & energy range\\\midrule
$10^{-4}$ (calibration default) & 0.39\,K & 13--24 & 37--62\\
$10^{-6}$ & -- & 0.36--0.64 & 1.0--2.1\\
$10^{-8}$ & \multicolumn{3}{c}{a block solve fails (both codes)}\\\bottomrule
\end{tabular}
\end{center}
At $10^{-6}$ the mean energies of V2.2 and this code agree to 0.11 (of 6945) at the prior centre; at a jittered point a difference of 0.55 at $10^{-6}$ shrinks to $0.035\pm0.036$ at $10^{-7}$, i.e.\ it is integration error, not a model difference. Trajectory and energy checks are therefore
made against this noise floor, not against fixed tolerances. For calibration the noise at $10^{-4}$ is of the order of the
energy differences the sampler uses; see Section~\ref{sec:gaps}.

The few V2.2 terms that were booked and then cancelled (lamp NIR on the canopy, V2.0 lamp heat on the air, floor
long-wave first taken from the air and then returned) are booked once here, so the results agree to rounding rather
than bit for bit.

\begin{center}\small
\begin{longtable}{>{\raggedright\arraybackslash}p{0.3\textwidth}>{\raggedright\arraybackslash}p{0.62\textwidth}}
\caption{Where each part of this document is in the code.}\\\toprule
document & code\\\midrule
state equations (\ref{eq:T1})--(\ref{eq:Tm}) & \code{model\_terms}, end of each block; \code{rhs!} (\code{src/climate.jl})\\
solar column & \code{solar\_column}, \code{column4}, \code{canopy\_layer}, \code{roof\_props} (\code{src/solar.jl})\\
plant light, FvCB & \code{PARp}, \code{assimilation}, \code{lai\_eff} (\code{src/photosynthesis.jl})\\
stomata & \code{q2}--\code{q10}, \code{q1}\\
cover node & \code{cover\_temperature}\\
long-wave & \code{r6}, \code{r7}, \code{r11}--\code{r14}\\
air exchange & \code{U8}, \code{f2}--\code{f7}\\
heating & \code{e}, \code{dmd}, \code{Tp}, \code{deint}, \code{h4}\\
lamps & \code{Ql}, \code{Qlc}, \code{Qla}\\
air heat, vapour, \ce{CO2} & \code{h1}--\code{h11}, \code{l1}, \code{l2}, \code{p1}--\code{p10}, \code{o1}--\code{oR}\\
floor, deep soil, mass & \code{Qaf}, \code{Qfd}, \code{Qfs}, \code{Qam}, \code{Qms}, \code{Cm}\\
inputs & \code{src/inputs.jl}, \code{src/agc\_data.jl}\\
crop sub-models & \code{node\_appearance}, \code{development\_rate}, \code{fruit\_sink}, \code{vegetative\_demand}, \code{cucumber\_generic} (\code{src/crop.jl})\\
LAI, daily step & \code{seasonal\_sla}, \code{crop\_step!}, \code{init\_crop\_state}\\
crop light, series & \code{daily\_assimilation\_measured}, \code{load\_measured\_climate\_par}, \code{crop\_series} (\code{src/crop\_measured.jl})\\
likelihoods & \code{climate\_energy} (\code{src/climate\_calibration.jl}), \code{crop\_energy}\\
\bottomrule
\end{longtable}
\end{center}


\end{document}
